\chapter{Benchmark Taxonomy Label Definitions}
\label{app:benchmark-taxonomy-labels}

This appendix defines the multi-label taxonomy used to annotate benchmark queries for later subgroup analysis. Each field captures a distinct property of the query or of the minimum evidence conditions required for a correct answer. 

\newcommand{\topcode}[1]{\parbox[t]{\linewidth}{\raggedright\ttfamily #1}}
\newcommand{\toptext}[1]{\parbox[t]{\linewidth}{\raggedright #1}}

\begingroup
\footnotesize
\setlength{\tabcolsep}{4pt}
\renewcommand{\arraystretch}{1.15}

\begin{longtable}{>{\raggedright\arraybackslash}p{2.6cm}
                  >{\raggedright\arraybackslash}p{3.3cm}
                  >{\raggedright\arraybackslash}p{3.6cm}
                  >{\raggedright\arraybackslash}p{4.8cm}}
\caption{Benchmark taxonomy label definitions and allowed values.}
\label{tab:benchmark-taxonomy-labels} \\

\hline
Field & What it captures & Allowed values & Annotation rule \\
\hline
\endfirsthead

\multicolumn{4}{l}{\textit{Table \thetable{} continued from previous page}} \\
\hline
Field & What it captures & Allowed values & Annotation rule \\
\hline
\endhead

\hline
\multicolumn{4}{r}{\textit{Continued on next page}} \\
\endfoot

\hline
\endlastfoot

\topcode{source\_\newline needed} &
\toptext{The minimum evidence class needed for a correct answer} &
\topcode{docs\newline tickets\newline both} &
\toptext{Labels describe the task, not the current retrieval result or the current vector-store contents. Assign the minimum evidence class required for a correct answer.} \\

\topcode{docs\_scope\_\newline needed} &
\toptext{Which class of official documentation is authoritative when documentary evidence is needed} &
\topcode{local\_official\newline external\_official\newline local\_and\_external\newline none} &
\toptext{Use \texttt{none} only when \texttt{source\_needed=tickets}. Queries with \texttt{source\_needed=docs} or \texttt{both} must not use \texttt{none}.} \\

\topcode{validity\_\newline sensitive} &
\toptext{Whether correctness depends on current system, partition, service, workflow, or compatibility truth} &
\topcode{yes\newline no} &
\toptext{Mark \texttt{yes} when current local operational truth is central to correctness; mark \texttt{no} when the answer is sufficiently stable.} \\

\topcode{attachment\_\newline dependent} &
\toptext{Whether visible text alone is sufficient, or whether unseen attachments materially determine correct handling} &
\topcode{yes\newline no} &
\toptext{Mark \texttt{yes} when screenshots, logs, uploaded documents, or other unseen attachments materially affect the correct first response.} \\

\topcode{query\_type} &
\toptext{The broad problem family represented by the query} &
\topcode{local\_operational\newline software\_version\newline general\_how\_to} &
\toptext{Assign the dominant problem family: local service truth, version-sensitive software or toolchain issues, or stable procedural guidance.} \\

\topcode{version\_\newline sensitive} &
\toptext{Whether correctness depends materially on a software or toolchain version} &
\topcode{yes\newline no} &
\toptext{Mark \texttt{yes} when version differences could change the correct guidance or resolution.} \\

\topcode{system\_scope\_\newline required} &
\toptext{Whether correctness depends on a particular system, service, partition, or platform} &
\topcode{yes\newline no} &
\toptext{Mark \texttt{yes} when the answer would differ materially across systems or execution environments.} \\

\end{longtable}
\endgroup

Where multiple interpretations were possible, labels were assigned according to the minimum evidence needed for a correct and locally applicable first-response answer. 